\scanpage{3}
Les autom. \inserted{à lacets} de $\hat\Pi_{0,3}=\pi$ commutant \deleted{\unsure{à}} $\sigma_0$, et commutant extérieurement à $\rho$, sont donnés par
\[
u(l_{\unsure{1|0}})=\operatorname{int}(g_1)l_1^{p}
\]
\[
u(l_\infty)=\operatorname{int}(\underbrace{\sigma_0(g_1)}_{g_\infty})l_\infty^{p}
\]
\margin{\uline{NB} on aura
\[
u(l_0)=\operatorname{int}\bigl(\underbrace{g_1l_1^{-\gamma}\rho^{-1}(h_1)}_{g_0}\bigr)\cdot l_0^{p}
\]}
avec $g_1\in\Pi\bmod L_1$, $p\in\hat{\mathbb Z}^{*}$ satisfaisant la condition (où $\beta,\gamma\in\hat{\mathbb Z}$ sont uniquement déterminés par $g_1,p$)
\margin{\uline{et $\bar m$ pas} \uline{$u$}!}
\[
(1)\qquad\boxed{l_\infty^{p}\cdot\operatorname{int}(h_1)(l_1^{p})\cdot\operatorname{int}\bigl(h_1l_1^{-\gamma}\rho^{-1}(h_1)\bigr)(l_0^{p})=1}
\]
et
\[
(2)\qquad\boxed{\rho(h_1)(l_\infty^{-\gamma}\unsure{\cdot|=} h_1l_1^{-\gamma})\rho^{-1}(h_1)=l_0^\beta}
\]
\margin{NB. On doit en fait avoir $\unsure{\beta=\gamma|\beta\ne\gamma}$ (passer à $(\Pi_{0,3})_{\mathrm{ab}}$)}
avec
\[
h_1=\unsure{\sigma_0|\sigma_0'|\dot\sigma_0}(g_1)^{-1}g_1
\]
La deuxième relation s'écrit aussi
\[
h_1l_1^{-\gamma}\rho^{-1}(h_1)=l_\infty^\gamma\rho(h_1)^{-1}l_0^{\deleted{\unsure{\gamma|\alpha}}\beta}
\]
et en portant dans la première relation, celle-ci peut se récrire
\figure{\begin{tikzpicture}[baseline]\draw[->] (0,0.7)--(0,0);\draw (-0.1,0.7)--(0.1,0.7);\end{tikzpicture}}
\[
(1')\qquad\boxed{l_\infty^{p}\cdot\unsure{\operatorname{int}(h_1)|\operatorname{int}h_1}(l_1^{p})\cdot\operatorname{int}\bigl(l_\infty^\gamma\rho(h_1)^{-1}\bigr)(l_0^{p})=1}
\]
\uline{NB} Si on change $g_1$ en $g_1'=g_1l_1^\mu$, $h_1$ est remplacé en $h_1'=l_\infty^{-\mu}h_1l_1^\mu$, et la relation $(1')$ donne, en lui appliquant $\operatorname{int}(l_\infty^{-\mu})$ :
\[
l_\infty^{p}\cdot\operatorname{int}(\deleted{\unsure{h_1}}h_1')(l_1^{p})\cdot\operatorname{int}\bigl(\deleted{\unsure{l_1^\mu|h_1}}l_\infty^{\unsure{\gamma'|\gamma-\mu|\gamma}}\unsure{\rho^{\deleted{\unsure{-1}}}|\rho^{-1}|\rho^1|\hat\rho}(h_1')^{-1}\bigr)(l_0^{p})=1
\]
(\uline{i.e.} \uline{$\gamma'=\gamma\deleted{\unsure{+2\mu|+\mu}}$}) et on a la relation correspondante \unsure{à|:} $(2)_2$ avec $\deleted{\unsure{\beta}}\beta'=\beta$.
\mspage{2}
